% content/indonesian/symbols.tex
% -----------------------------------------------------------------------------
% DAFTAR SIMBOL
% -----------------------------------------------------------------------------
% Muat hanya notasi inti yang dipakai lintas bab, dengan deskripsi satu baris
% agar daftar muat satu halaman. Simbol yang hanya muncul pada satu persamaan
% cukup didefinisikan langsung di bawah persamaannya.
%
% Urutan: huruf Latin, huruf Yunani, lalu simbol kaligrafi. Tuliskan satuan
% bila ada. Gunakan notasi yang persis sama dengan yang dipakai pada bab.
%
% Judul halaman ditulis oleh content/frontmatter-lists.tex, jadi berkas ini
% cukup berisi tabelnya saja.

\TemplateTodo{Ganti entri contoh di bawah ini dengan notasi yang benar-benar dipakai pada naskah Anda. Jika naskah tidak memuat persamaan, hapus blok DAFTAR SIMBOL pada content/frontmatter-lists.tex.}

\begin{longtable}{@{}p{0.15\linewidth}p{0.80\linewidth}@{}}
$x_i$ & Masukan sampel ke-$i$. \\
$y_i$, $\hat{y}_i$ & Nilai sebenarnya dan nilai prediksi sampel ke-$i$. \\
$N$ & Jumlah sampel pada himpunan yang dievaluasi. \\
$\eta$ & Laju pembelajaran (\textit{learning rate}). \\
$\theta$ & Seluruh parameter model yang dilatih. \\
$\mathcal{L}$ & Fungsi \textit{loss} yang diminimalkan selama pelatihan. \\
\end{longtable}
